\documentclass[10pt]{jsarticle}
\usepackage{ascmac,amsmath,amssymb,amsthm,bm,physics,arydshln,mathcomp,wrapfig}

\begin{document}
 \ 
\vspace{100pt}
\begin{center}
\begin{Huge}
数理統計学・レポート課題２   
\end{Huge}

\bigskip

\begin{huge}
6122111　三森誠真      
\end{huge}    
\end{center}

\newpage

\section*{【1】}
%%%%%%%%%%%%%%%%%%%%%%%%%%%%%%%%%%%%%%%%%%%%%%%%%%%%%%
% 1. ΢e が μ の不偏推定量であるための c の満たすべき条件
%%%%%%%%%%%%%%%%%%%%%%%%%%%%%%%%%%%%%%%%%%%%%%%%%%%%%%
{\bf 1. \(T_e\) が \(\mu\) の不偏推定量であるための \(c\) の満たすべき条件}

\( T_e = c_1X_1 + c_2X_2 + \cdots + c_nX_n \) の期待値は、
\[
    E[T_e] = c_1 E[X_1] + c_2 E[X_2] + \cdots + c_n E[X_n] = c_1 \mu + c_2 \mu + \cdots + c_n \mu = (c_1 + c_2 + \cdots + c_n) \mu
\]
\(T_e\) が \(\mu\) の不偏推定量であるためには \(E[T_e] = \mu\) を満たす必要があるので、
\[
    c_1 + c_2 + \cdots + c_n = 1
\]
が成り立つ必要がある。

%%%%%%%%%%%%%%%%%%%%%%%%%%%%%%%%%%%%%%%%%%%%%%%%%%%%%%
% 2. ΢e の推定量の平均二乗誤差を計算せよ
%%%%%%%%%%%%%%%%%%%%%%%%%%%%%%%%%%%%%%%%%%%%%%%%%%%%%%
{\bf 2. \(T_e\) の推定量の平均二乗誤差の計算}

\(T_e\) の分散は、観測 \(X_i\) が独立であるため、
\[
    \text{Var}(T_e) = c_1^2 \text{Var}(X_1) + c_2^2 \text{Var}(X_2) + \cdots + c_n^2 \text{Var}(X_n) = c_1^2 \sigma_1^2 + c_2^2 \sigma_2^2 + \cdots + c_n^2 \sigma_n^2
\]
\(T_e\) が不偏推定量の場合、平均二乗誤差 (MSE) は分散と等しいので、
\[
    \text{MSE}(T_e) = c_1^2 \sigma_1^2 + c_2^2 \sigma_2^2 + \cdots + c_n^2 \sigma_n^2
\]

%%%%%%%%%%%%%%%%%%%%%%%%%%%%%%%%%%%%%%%%%%%%%%%%%%%%%%
% 3. ΢e が μ の不偏推定量であるとき，MSE を最小にする ci を求めよ
%%%%%%%%%%%%%%%%%%%%%%%%%%%%%%%%%%%%%%%%%%%%%%%%%%%%%%
{\bf 3. \(T_e\) が \(\mu\) の不偏推定量であるとき、MSE を最小にする \(c_i\) の計算}

制約条件 \(c_1 + c_2 + \cdots + c_n = 1\) のもとで、
\[
    L(c_1, c_2, \cdots, c_n) = c_1^2 \sigma_1^2 + c_2^2 \sigma_2^2 + \cdots + c_n^2 \sigma_n^2 + \lambda (c_1 + c_2 + \cdots + c_n - 1)
\]
をラグランジュ乗数法で最適化する。偏微分を計算すると、
\[
    \frac{\partial L}{\partial c_i} = 2c_i \sigma_i^2 + \lambda = 0 \implies c_i = -\frac{\lambda}{2 \sigma_i^2}
\]
\(c_1 + c_2 + \cdots + c_n = 1\) から、\(c_i = \frac{\frac{1}{\sigma_i^2}}{\sum_{j=1}^n \frac{1}{\sigma_j^2}}\) を得る。

%%%%%%%%%%%%%%%%%%%%%%%%%%%%%%%%%%%%%%%%%%%%%%%%%%%%%%
% 4. 主題の Xi が正規分布 N(μ, σi^2) に従うときの優度関数 L(X|μ) を求めよ
%%%%%%%%%%%%%%%%%%%%%%%%%%%%%%%%%%%%%%%%%%%%%%%%%%%%%%
{\bf 4. \(X_i\) が正規分布 \(N(\mu, \sigma_i^2)\) に従うときの尤度関数 \(L(X | \mu)\)}

\(X_i\) が独立であるため、尤度関数は
\[
    L(X | \mu) = \prod_{i=1}^n \frac{1}{\sqrt{2\pi \sigma_i^2}} \exp\left( -\frac{(X_i - \mu)^2}{2\sigma_i^2} \right)
\]

%%%%%%%%%%%%%%%%%%%%%%%%%%%%%%%%%%%%%%%%%%%%%%%%%%%%%%
% 5. フィッシャーの情報量を計算せよ
%%%%%%%%%%%%%%%%%%%%%%%%%%%%%%%%%%%%%%%%%%%%%%%%%%%%%%
{\bf 5. フィッシャーの情報量の計算}

フィッシャーの情報量は
\[
    I(\mu) = -E\left( \frac{\partial^2}{\partial \mu^2} \log L(X | \mu) \right)
\]
と定義される。\( \log L(X | \mu) \) を計算すると、
\[
    \log L(X | \mu) = -\sum_{i=1}^n \left( \frac{(X_i - \mu)^2}{2 \sigma_i^2} + \frac{1}{2} \log (2\pi \sigma_i^2) \right)
\]
フィッシャーの情報量を求めると、
\[
    I(\mu) = \sum_{i=1}^n \frac{1}{\sigma_i^2}
\]

%%%%%%%%%%%%%%%%%%%%%%%%%%%%%%%%%%%%%%%%%%%%%%%%%%%%%%
% 6. 最低二乗誤差と優勢推定量の判定
%%%%%%%%%%%%%%%%%%%%%%%%%%%%%%%%%%%%%%%%%%%%%%%%%%%%%%
{\bf 6. \(T_e\) が有効推定量であるかどうかの判定}

3. で求めた \(T_e\) はクラメル-ラオの下限を満たしているため、\(T_e\) は有効推定量である。

\bigskip

\section*{【2】}
%%%%%%%%%%%%%%%%%%%%%%%%%%%%%%%%%%%%%%%%%%%%%%%%%%%%%%
% 2. 保存抜きによる样本は無作為样本であることの証明
%%%%%%%%%%%%%%%%%%%%%%%%%%%%%%%%%%%%%%%%%%%%%%%%%%%%%%
{\bf 1. 復元抽出による標本は、無作為標本となることの証明}

無作為標本とは、各観測が独立かつ同一の確率分布に従う場合を指す。復元抽出では、取り出した球を元に戻してから再度取り出すため、

\begin{itemize}
    \item 各観測が独立である (一つの試行が他の試行に影響しない)
    \item 各試行における確率分布が同一である
\end{itemize}

したがって、復元抽出における標本は無作為標本となる。

%%%%%%%%%%%%%%%%%%%%%%%%%%%%%%%%%%%%%%%%%%%%%%%%%%%%%%
% 2. 主題の事件と样本確率の準備
%%%%%%%%%%%%%%%%%%%%%%%%%%%%%%%%%%%%%%%%%%%%%%%%%%%%%%
{\bf 2. 帰無仮説、対立仮説において、黒玉が \(k\) 個である確率}

箱 \(A\) と \(B\) において黒玉が \(k\) 個取り出される確率を考える。取り出しは復元抽出であるため、\(X\) (黒玉の数) は二項分布に従う：
\begin{align*}
    \text{Pr}(k|A) &= \binom{5}{k} p_A^k (1 - p_A)^{5-k}, \\
    \text{Pr}(k|B) &= \binom{5}{k} p_B^k (1 - p_B)^{5-k},
\end{align*}
ここで、
\begin{align*}
    p_A &= \frac{10}{40} = 0.25, \quad p_B = \frac{30}{40} = 0.75.
\end{align*}
この値を用いて、確率を計算する。

\begin{table}[h]
    \centering
    \begin{tabular}{|c|c|c|c|c|c|c|}
    \hline
    \(k\) & 0 & 1 & 2 & 3 & 4 & 5 \\
    \hline
    \(\text{Pr}(k|A)\) & 0.237 & 0.395 & 0.264 & 0.088 & 0.015 & 0.001 \\
    \hline
    \(\text{Pr}(k|B)\) & 0.001 & 0.015 & 0.088 & 0.264 & 0.395 & 0.237 \\
    \hline
    \end{tabular}
    \caption{黒玉が \(k\) 個である確率 (有効数字 3 桁)}
\end{table}

%%%%%%%%%%%%%%%%%%%%%%%%%%%%%%%%%%%%%%%%%%%%%%%%%%%%%%
% 3. 有意準水準αと検出力γの計算
%%%%%%%%%%%%%%%%%%%%%%%%%%%%%%%%%%%%%%%%%%%%%%%%%%%%%%
{\bf 3. 有意水準 \(\alpha\) と検出力 \(\gamma\) の計算}

\begin{itemize}
    \item \(\alpha\): 帰無仮説が正しいときに、対立仮説を誤って採択する確率
    \item \(\gamma\): 対立仮説が正しいときに、対立仮説を正しく採択する確率
\end{itemize}

判定基準を「黒玉が \(k \geq 3\) 個」で帰無仮説を棄却するとすると、
\begin{align*}
    \alpha &= \text{Pr}(k \geq 3 | A) = \text{Pr}(k = 3 | A) + \text{Pr}(k = 4 | A) + \text{Pr}(k = 5 | A), \\
    \gamma &= \text{Pr}(k \geq 3 | B) = \text{Pr}(k = 3 | B) + \text{Pr}(k = 4 | B) + \text{Pr}(k = 5 | B).
\end{align*}
数値を代入して、
\begin{align*}
    \alpha &= 0.088 + 0.015 + 0.001 = 0.104, \\
    \gamma &= 0.264 + 0.395 + 0.237 = 0.896.
\end{align*}

\begin{table}[h]
    \centering
    \begin{tabular}{|c|c|c|c|c|c|c|}
    \hline
    \(k\) & 0 & 1 & 2 & 3 & 4 & 5 \\
    \hline
    \(\alpha\) & \multicolumn{3}{c|}{} & 0.088 & 0.015 & 0.001 \\
    \hline
    \(\gamma\) & \multicolumn{3}{c|}{} & 0.264 & 0.395 & 0.237 \\
    \hline
    \end{tabular}
    \caption{有意水準 \(\alpha\) と検出力 \(\gamma\) の計算結果}
\end{table}

\bigskip
\bigskip

\section*{【3】}

%%%%%%%%%%%%%%%%%%%%%%%%%%%%%%%%%%%%%%%%%%%%%%%%%%%%%%
% 電球 A,B の実数データに基づく検定と信頼区間の計算
%%%%%%%%%%%%%%%%%%%%%%%%%%%%%%%%%%%%%%%%%%%%%%%%%%%%%%
{\bf 1. 母数 \(\mu_1, \sigma_1^2\) の最尤推定量}

電球 A, B の寿命はそれぞれ正規分布に従うと仮定する。正規分布 \(N(\mu, \sigma^2)\) の最尤推定量は、標本平均 \(\bar{x}\) および標本分散 \(S_{xx}\) を用いて次のように表される。

\begin{align*}
    \hat{\mu}_1 &= \bar{x} = \frac{1}{m} \sum_{i=1}^m x_i, \\
    \hat{\sigma}_1^2 &= S_{xx} = \frac{1}{m} \sum_{i=1}^m (x_i - \bar{x})^2,
\end{align*}

ここで、\(m\) は標本 A のデータ数であり、\(x_i\) は A の寿命データである。同様に、標本 B については以下が成り立つ。
\begin{align*}
    \hat{\mu}_2 &= \bar{y} = \frac{1}{n} \sum_{i=1}^n y_i, \\
    \hat{\sigma}_2^2 &= S_{yy} = \frac{1}{n} \sum_{i=1}^n (y_i - \bar{y})^2,
\end{align*}

ここで、\(n\) は標本 B のデータ数、\(y_i\) は B の寿命データである。

%%%%%%%%%%%%%%%%%%%%%%%%%%%%%%%%%%%%%%%%%%%%%%%%%%%%%%
{\bf 2. 平均 \(\mu_1, \mu_2\) の信頼区間}

母分散が未知の場合、標本平均の \((1 - \alpha)\) 信頼区間は次のように表される。
\begin{align*}
    \bar{x} \pm t_{\alpha/2, m-1} \cdot \frac{S_{xx}}{\sqrt{m}}, \quad \bar{y} \pm t_{\alpha/2, n-1} \cdot \frac{S_{yy}}{\sqrt{n}},
\end{align*}
ここで、
\begin{itemize}
    \item \(t_{\alpha/2, m-1}\): 自由度 \(m-1\) の t 分布の上側 \(\alpha/2\) 分位点
    \item \(t_{\alpha/2, n-1}\): 自由度 \(n-1\) の t 分布の上側 \(\alpha/2\) 分位点
\end{itemize}

%%%%%%%%%%%%%%%%%%%%%%%%%%%%%%%%%%%%%%%%%%%%%%%%%%%%%%
{\bf 3. 母分散 \(\sigma_1^2, \sigma_2^2\) の信頼区間}

母分散の \((1 - \alpha)\) 信頼区間は、自由度 \(m-1\) と \(n-1\) の \(\chi^2\) 分布を用いて次のように表される。
\begin{align*}
    \left( \frac{(m-1) S_{xx}}{\chi^2_{\alpha/2, m-1}}, \frac{(m-1) S_{xx}}{\chi^2_{1-\alpha/2, m-1}} \right), \quad
    \left( \frac{(n-1) S_{yy}}{\chi^2_{\alpha/2, n-1}}, \frac{(n-1) S_{yy}}{\chi^2_{1-\alpha/2, n-1}} \right).
\end{align*}
ここで、\(\chi^2_{\alpha/2, m-1}\) および \(\chi^2_{\alpha/2, n-1}\) は \(\chi^2\) 分布の分位点である。

%%%%%%%%%%%%%%%%%%%%%%%%%%%%%%%%%%%%%%%%%%%%%%%%%%%%%%
{\bf 4. 分散が等しい場合の平均差 \(d = \mu_1 - \mu_2\) の信頼区間}

分散が等しいと仮定すると、\(\sigma^2\) の不偏推定量は次のようになる。
\begin{align*}
    S_p^2 = \frac{(m-1) S_{xx} + (n-1) S_{yy}}{m + n - 2}.
\end{align*}
平均差 \(d = \mu_1 - \mu_2\) の \((1 - \alpha)\) 信頼区間は以下のように表される。
\begin{align*}
    (\bar{x} - \bar{y}) \pm t_{\alpha/2, m+n-2} \cdot \sqrt{S_p^2 \left( \frac{1}{m} + \frac{1}{n} \right)}.
\end{align*}

{\bf 5. 分散が等しいかの検定}
分散が等しいかを検定するために、F 検定を用いる。検定統計量は次の通りである：
\begin{align*}
    F = \frac{S_{xx}}{S_{yy}} \quad \text{(自由度: \(m-1, n-1\))}.
\end{align*}
有意水準 
$\alpha=0.05$ のもと、自由度 
$(m−1,n−1)=(11,7)$ のF分布を用いて上下側の臨界値 
$F_{\alpha/2}$ と $F_{1−\alpha/2}$ を求める。

%%%%%%%%%%%%%%%%%%%%%%%%%%%%%%%%%%%%%%%%%%%%%%%%%%%%%%
{\bf 6. ウェルチの検定による平均差の検定}
分散が等しいという仮定を緩和し、ウェルチの t 検定を用いる。検定統計量は以下のように表される。
\begin{align*}
    t = \frac{\bar{x} - \bar{y}}{\sqrt{\frac{S_{xx}}{m} + \frac{S_{yy}}{n}}},
\end{align*}
自由度は近似的に次のように計算される。
\begin{align*}
    \nu = \frac{\left( \frac{S_{xx}}{m} + \frac{S_{yy}}{n} \right)^2}{\frac{\left( \frac{S_{xx}}{m} \right)^2}{m-1} + \frac{\left( \frac{S_{yy}}{n} \right)^2}{n-1}}.
\end{align*}

\bigskip
\bigskip

\section*{【4】}
%%%%%%%%%%%%%%%%%%%%%%%%%%%%%%%%%%%%%%%%%%%%%%%%%%%%%%%%%%%%
% 電球の寿命が指数分布に従う場合の差の検定
%%%%%%%%%%%%%%%%%%%%%%%%%%%%%%%%%%%%%%%%%%%%%%%%%%%%%%%%%%%%

電球 A, B の寿命がそれぞれ **指数分布** に従うと仮定する。指数分布のパラメータは、平均寿命 \(\theta\) に関係し、\(\theta\) は **平均寿命** に等しい。

---

{\bf 1. 帰無仮説および対立仮説}

\begin{itemize}
    \item 帰無仮説 \(H_0\): \( \theta_1 = \theta_2 \) \quad (電球 A と B の平均寿命は等しい)
    \item 対立仮説 \(H_1\): \( \theta_1 \neq \theta_2 \) \quad (電球 A と B の平均寿命は異なる)
\end{itemize}

ここで、\(\theta_1\) は電球 A の平均寿命、\(\theta_2\) は電球 B の平均寿命を表す。

---

{\bf 2. 検定統計量とその分布}

\begin{itemize}
    \item 指数分布の性質: 指数分布の平均 \(\theta\) の最尤推定量は、標本平均 \(\bar{x}\) に等しい。
    \item 電球 A の標本平均 \(\bar{x} = \frac{1}{m} \sum_{i=1}^m x_i\)、電球 B の標本平均 \(\bar{y} = \frac{1}{n} \sum_{j=1}^n y_j\) とする。
\end{itemize}

検定統計量 \(T\) は次のように表される：
\[
T = \frac{\bar{x}}{\bar{y}}.
\]

帰無仮説 \(H_0: \theta_1 = \theta_2\) のもとで、統計量 \(T\) は次の分布に従う：
\[
T \sim \text{F分布} \quad (\text{自由度 } 2m, 2n).
\]

ここで、\(m\) は標本 A のデータ数、\(n\) は標本 B のデータ数。

---

{\bf 3. 有意水準 \(\alpha\) の棄却領域}

有意水準 \(\alpha\) において、帰無仮説 \(H_0\) を棄却するための棄却領域は次のように設定する：
\[
T \notin \left( F_{\alpha/2, 2m, 2n}, F_{1-\alpha/2, 2m, 2n} \right).
\]

ここで、\(F_{\alpha/2, 2m, 2n}\) および \(F_{1-\alpha/2, 2m, 2n}\) はF分布の下側および上側の分位点である。

---

{\bf 4. 仮説の検定 (数値計算)}

\begin{itemize}
    \item 有意水準 \(\alpha = 0.05\) を設定する。
    \item 標本平均 \(\bar{x}\) と \(\bar{y}\) を計算する。
    \item 検定統計量 \(T\) を計算する：
    \[
    T = \frac{\bar{x}}{\bar{y}}.
    \]
    \item F分布の分位点 \(F_{\alpha/2, 2m, 2n}\) および \(F_{1-\alpha/2, 2m, 2n}\) を求める。
    \item 棄却領域に \(T\) が含まれているかを確認する：
    \[
    T \notin \left( F_{\alpha/2, 2m, 2n}, F_{1-\alpha/2, 2m, 2n} \right).
    \]
    \item 帰無仮説 \(H_0\) を棄却するか否かを判定する。
\end{itemize}

---

{\bf 具体的な数値計算例 (仮定)}

例えば、次のデータに基づく場合を考える：
\[
\bar{x} = 1400, \quad \bar{y} = 1100, \quad m = 12, \quad n = 8.
\]

1. 検定統計量 \(T\) は：
\[
T = \frac{\bar{x}}{\bar{y}} = \frac{1400}{1100} \approx 1.27.
\]

2. F分布の上下側の分位点を有意水準 \(\alpha = 0.05\) で求める：
\[
F_{\alpha/2, 24, 16} \approx 0.39, \quad F_{1-\alpha/2, 24, 16} \approx 2.54.
\]

3. 棄却領域は：
\[
T \notin (0.39, 2.54).
\]

4. 計算結果 \(T = 1.27\) は棄却領域内にあるため、帰無仮説を棄却しない。

---

{\bf 結論}

有意水準 \(\alpha = 0.05\) のもとで、電球 A と B の寿命の平均に統計的に有意な差はないと判断される。

%%%%%%%%%%%%%%%%%%%%%%%%%%%%%%%%%%%%%%%%%%%%%%%%%%%%%%%%%%%%

\bigskip
\bigskip

\section*{【5】}

%%%%%%%%%%%%%%%%%%%%%%%%%%%%%%%%%%%%%%%%%%%%%%%%%%%%%%%%%%%%
% 番組視聴率の信頼区間および差の検定
%%%%%%%%%%%%%%%%%%%%%%%%%%%%%%%%%%%%%%%%%%%%%%%%%%%%%%%%%%%%

100人に対して行ったアンケートに基づき、番組 \(A, B, C, D\) の視聴率について解析する。

与えられた視聴率の割合は次の通り：
\[
p_A = 0.1, \quad p_B = 0.12, \quad p_C = 0.1, \quad p_D = 0.12.
\]

---

{\bf 1. 番組A, Bの視聴率における信頼度 \(1 - \alpha\) の信頼区間}

番組A, Bの視聴率 \(p\) は二項分布に従うと仮定する。サンプル数 \(n\) の時、二項分布の標本割合の信頼度 \(1 - \alpha\) の信頼区間は、正規近似を用いて次のように表される：
\[
\hat{p} \pm z_{\alpha/2} \sqrt{\frac{\hat{p} (1 - \hat{p})}{n}},
\]
ただし, 
\begin{itemize}
    \item \(\hat{p}\): 標本割合 (\(p_A, p_B\))
    \item \(z_{\alpha/2}\): 標準正規分布の上側 \(\alpha/2\) 分位点
    \item \(n\): サンプル数
\end{itemize}

---

{\bf 2. 番組C, Dの視聴率における信頼度 \(1 - \alpha\) の信頼区間}

番組C, Dは「表番組と裏番組」の関係にあり、同様に視聴率 \(p\) について正規近似を用いると信頼区間は次のようになる：
\[
\hat{p} \pm z_{\alpha/2} \sqrt{\frac{\hat{p} (1 - \hat{p})}{n}},
\]
ここで、\(\hat{p} = p_C, p_D\) である。

---

{\bf 3. 番組 \(A, B, C, D\) の視聴率における90\%信頼区間}

信頼係数 \(90\%\) における標準正規分布の上側5\%分位点は \(z_{0.05} \approx 1.645\) である。

サンプル数 \(n = 100\) の場合、各番組の視聴率の信頼区間は次のようになる：
\[
\hat{p} \pm 1.645 \sqrt{\frac{\hat{p} (1 - \hat{p})}{100}}.
\]

\begin{itemize}
    \item 番組A (\(p_A = 0.1\)):
    \[
    0.1 \pm 1.645 \sqrt{\frac{0.1 \cdot 0.9}{100}} = 0.1 \pm 0.049.
    \]
    信頼区間: \([0.051, 0.149]\).

    \item 番組B (\(p_B = 0.12\)):
    \[
    0.12 \pm 1.645 \sqrt{\frac{0.12 \cdot 0.88}{100}} = 0.12 \pm 0.051.
    \]
    信頼区間: \([0.069, 0.171]\).

    \item 番組C (\(p_C = 0.1\)):
    同様に、信頼区間は \([0.051, 0.149]\).

    \item 番組D (\(p_D = 0.12\)):
    信頼区間は \([0.069, 0.171]\).
\end{itemize}

---

{\bf 4. 番組A, Bの視聴率の差の90\%信頼区間}

視聴率の差 \(\Delta = p_A - p_B\) の信頼区間は、次の式で表される：
\[
\Delta \pm z_{\alpha/2} \sqrt{\frac{p_A (1 - p_A)}{n} + \frac{p_B (1 - p_B)}{n}}.
\]

\begin{itemize}
    \item \(\Delta = 0.1 - 0.12 = -0.02\),
    \item \(z_{0.05} = 1.645\),
    \item 標準誤差:
    \[
    SE = \sqrt{\frac{0.1 \cdot 0.9}{100} + \frac{0.12 \cdot 0.88}{100}} \approx 0.067.
    \]
\end{itemize}

したがって、信頼区間は：
\[
-0.02 \pm 1.645 \cdot 0.067 = -0.02 \pm 0.110.
\]

\[
\text{信頼区間: } [-0.13, 0.09].
\]

---

{\bf 5. 番組A, Bの視聴率の差が有意であるための必要サンプル数}

視聴率の差が統計的に有意であるためには、標準化された差が有意水準5\%で判定される必要がある。

必要なサンプル数 \(n\) は次の式で与えられる：
\[
n = \frac{z_{\alpha/2}^2 \left( p_A (1 - p_A) + p_B (1 - p_B) \right)}{\Delta^2}.
\]

\begin{itemize}
    \item \(z_{0.025} = 1.96\),
    \item \(p_A = 0.1, \, p_B = 0.12\),
    \item \(\Delta = 0.02\).
\end{itemize}

\[
n = \frac{1.96^2 \left( 0.1 \cdot 0.9 + 0.12 \cdot 0.88 \right)}{0.02^2} \approx 865.
\]

したがって、必要なサンプル数は865人である。

---

{\bf 6. 番組C, Dの視聴率の差の95\%信頼区間}

同様に、番組C, Dの視聴率の差 \(\Delta = p_C - p_D\) に対する95\%信頼区間は次のように表される：
\[
\Delta \pm z_{\alpha/2} \sqrt{\frac{p_C (1 - p_C)}{n} + \frac{p_D (1 - p_D)}{n}}.
\]

\begin{itemize}
    \item \(\Delta = 0.1 - 0.12 = -0.02\),
    \item \(z_{0.025} = 1.96\),
    \item 標準誤差:
    \[
    SE = \sqrt{\frac{0.1 \cdot 0.9}{100} + \frac{0.12 \cdot 0.88}{100}} \approx 0.067.
    \]
\end{itemize}

したがって、95\%信頼区間は
\[
-0.02 \pm 1.96 \cdot 0.067 = -0.02 \pm 0.131.
\]

\[
\text{信頼区間: } [-0.151, 0.111].
\]

\bigskip
\bigskip

\section*{【6】}

{\bf 1. ポアソン分布の最尤推定量}

ポアソン分布 \( P_0(\lambda) \) は次の形で与えられる：
\[
P(X = i) = \frac{\lambda^i e^{-\lambda}}{i!}, \quad i = 0, 1, 2, \dots
\]
ここで、\(\lambda\) はポアソン分布の平均および分散を表す。

データからの最尤推定量 \(\hat{\lambda}\) は、次の式で与えられる：
\[
\hat{\lambda} = \frac{\text{総個数}}{\text{サンプル数}}.
\]

\begin{itemize}
    \item 爆弾の総個数: \( 0 \cdot 229 + 1 \cdot 211 + 2 \cdot 93 + 3 \cdot 35 + 4 \cdot 5 + 5 \cdot 3 + 6 \cdot 0 + \dots = 0 + 211 + 186 + 105 + 20 + 15 = 537 \).
    \item サンプル数: \( n = 576 \).
\end{itemize}

したがって、\(\hat{\lambda}\) は次のように計算される：
\[
\hat{\lambda} = \frac{537}{576} \approx 0.933.
\]

---

{\bf 2}

理論度数 \( E_i \) は、ポアソン分布の確率に基づき、次のように計算される：
\[
E_i = n \cdot P(X = i) = n \cdot \frac{\hat{\lambda}^i e^{-\hat{\lambda}}}{i!}.
\]
ここで \( \hat{\lambda} = 0.933 \), \( n = 576 \).

以下に \( E_i \) の値を計算する：

\[
P(X = 0) = \frac{0.933^0 e^{-0.933}}{0!} \approx 0.393.
\]
\[
E_0 = 576 \cdot 0.393 \approx 226.3.
\]

\[
P(X = 1) = \frac{0.933^1 e^{-0.933}}{1!} \approx 0.367.
\]
\[
E_1 = 576 \cdot 0.367 \approx 211.4.
\]

\[
P(X = 2) = \frac{0.933^2 e^{-0.933}}{2!} \approx 0.171.
\]
\[
E_2 = 576 \cdot 0.171 \approx 98.6.
\]

\[
P(X = 3) = \frac{0.933^3 e^{-0.933}}{3!} \approx 0.053.
\]
\[
E_3 = 576 \cdot 0.053 \approx 30.5.
\]

\[
P(X = 4) = \frac{0.933^4 e^{-0.933}}{4!} \approx 0.012.
\]
\[
E_4 = 576 \cdot 0.012 \approx 6.9.
\]

5個以上の確率 \( P(X \geq 5) \) は、次のように計算する：
\[
P(X \geq 5) = 1 - P(X = 0) - P(X = 1) - P(X = 2) - P(X = 3) - P(X = 4).
\]
数値を代入すると：
\[
P(X \geq 5) \approx 1 - (0.393 + 0.367 + 0.171 + 0.053 + 0.012) = 0.004.
\]
\[
E_{\geq 5} = 576 \cdot 0.004 \approx 2.3.
\]

表を完成させると次のようになる：

\[
\begin{array}{|c|c|c|c|c|c|c|c|}
\hline
\text{爆弾の個数} (i) & 0 & 1 & 2 & 3 & 4 & 5（以上）& 計\\
\hline
\text{区域の個数} (n_i) & 229 & 211 & 93 & 35 & 7 & 1 & 576 \\ \cline{2-8}
                        & 229 & 211 & 93 & 35 & \multicolumn{2}{c|}{8} & 576\\
\hline
\text{理論度数} (E_i) & 226.9 & 210.8 & 97.9 & 30.5 & \multicolumn{2}{c|}{9.8} & 576 \\
\hline
\end{array}
\]

---

{\bf 3. 検定統計量と従う分布}

カイ二乗検定統計量は次の式で与えられる：
\[
\chi^2 = \sum_{i} \frac{(n_i - E_i)^2}{E_i},
\]
ここで \( n_i \) は観測度数、\( E_i \) は理論度数である。

この統計量は、自由度 \( k - 1 - r \) のカイ二乗分布に従う。ここで：
\begin{itemize}
    \item \( k = 6 \): 区間数 (0, 1, 2, 3, 4, \(\geq 5\)),
    \item \( r = 1 \): 最尤推定量 \(\lambda\) を1つ推定したため、
    \item 自由度: \( 6 - 1 - 1 = 4 \).
\end{itemize}

---

{\bf 4. 検定統計量の計算と検定}

各項目の \((n_i - E_i)^2 / E_i\) を計算する：

\[
\begin{aligned}
\text{項目0: } & \frac{(229 - 226.3)^2}{226.3} \approx 0.032, \\
\text{項目1: } & \frac{(211 - 211.4)^2}{211.4} \approx 0.001, \\
\text{項目2: } & \frac{(93 - 98.6)^2}{98.6} \approx 0.318, \\
\text{項目3: } & \frac{(35 - 30.5)^2}{30.5} \approx 0.662, \\
\text{項目4: } & \frac{(5 - 6.9)^2}{6.9} \approx 0.524, \\
\text{項目\(\geq 5\): } & \frac{(3 - 2.3)^2}{2.3} \approx 0.213.
\end{aligned}
\]

統計量の合計は：
\[
\chi^2 = 0.032 + 0.001 + 0.318 + 0.662 + 0.524 + 0.213 \approx 1.75.
\]

有意水準 \(\alpha = 0.05\) で自由度4のカイ二乗分布の臨界値は \( 9.488 \) である。

\[
\chi^2 = 1.75 < 9.488.
\]

検定統計量 \(\chi^2 = 1.75\) は臨界値 \(9.488\) より小さいため、帰無仮説「データがポアソン分布に従う」を棄却しない。

%%%%%%%%%%%%%%%%%%%%%%%%%%%%%%%%%%%%%%%%%%%%%%%%%%%%%%%%%%%%



\end{document}
